綜合第六章之結論，本研究透過兩大訓練範式的系統性跨資料集比較，揭示了當前深偽偵測系統普遍面臨的跨域泛化瓶頸，並透過消融實驗確認了 Prompt Tuning 範式受制於前端硬性離散化的量化資訊損失，以及預訓練—微調範式受限於後端分類器特徵融合能力不足的結構性問題。這些發現指出，未來的深偽防禦系統不應僅追求單一資料集上的準確率極大化，更應著眼於縮小兩大範式之間的效能差距、提升跨域泛化能力，以及探索更具環境強健性的模型架構。本章針對上述方向提出四項未來展望。

基於第六章結論所揭示的離散化瓶頸，強制特徵離散化會導致用以辨別真偽的微小瑕疵（如高頻假影或壓縮損失）被當作雜訊過濾；然而，W2V-AASIST 所代表的連續聲學特徵模型雖然精準，卻可能在面對具有極端變異（如 MLAAD 資料集）的場景時，反而受到與深偽本質無關的聲學變異所誤導。最直接的改進路徑是構建「雙流架構 (Two-stream Architecture)」，同時輸入未經壓縮的連續波形特徵與具有高階語境意義的離散 Token，將兩種範式的優勢加以融合。這樣的設計能讓模型既兼顧局部的聲學瑕疵偵測，又能捕捉全域的序列與語意異常，有望在不犧牲 Prompt Tuning 範式訓練效率的前提下，彌補當前離散化架構在細粒度聲學特徵上的先天不足。

更根本的解決方向是探索原生支援連續聲學特徵輸入的新一代語音語言模型。本研究的可學習群集中心實驗已初步驗證，透過線性投影取代硬性 K-Means 量化可將 EER 由 18.50\% 大幅降至 3.84\%，但此初步探索仍存在數項待解之議題。首先，本研究已完成跨域驗證，結果顯示直接線性映射的改善並非對同域資料集的過擬合——其在跨域測試集上同樣帶來一致的效能躍升（ASV21 LA 由 18.23\% 降至 9.08\%，ASV21 DF 由 28.94\% 降至 14.44\%）。然而，在最具挑戰性的 ASV21 DF 上，直接線性映射（14.44\%）與 W2V-AASIST（5.57\%）之間仍存在顯著差距，且差距隨測試集難度增加而擴大，顯示即使消除前端量化瓶頸後，GSLM 骨幹架構的離散化推理機制仍是制約跨域泛化的深層限制。其次，本研究所採用的 GSLM 語言模型在設計上仍以離散 Token 序列為輸入前提，將連續特徵強行投影至其輸入空間，本質上並未改變語言模型內部的離散化推理機制。近年來，以擴散模型（Diffusion Model）或流匹配（Flow Matching）技術為基礎的新一代語音生成模型已展示了直接在連續聲學特徵空間中進行生成建模的能力——例如 NaturalSpeech~2 採用潛在擴散模型對連續編解碼器潛在向量進行建模~\cite{shen2024naturalspeech2}，Voicebox 則以流匹配框架直接對連續梅爾頻譜特徵進行推論~\cite{le2023voicebox}——此類架構為語音語言模型擺脫離散 Token 的前提限制提供了具體可行的實現路徑。未來若能以此類原生支援連續輸入的語音語言模型替換 GSLM 骨幹，使 Prompt Tuning 範式在保有大語言模型泛化推理能力的同時從根本上解決量化資訊損失問題，兩種範式各自的優勢將有望在同一框架內整合。

在跨域強健性方面，目前的模型在面對 InTheWild 與 ASV21 DeepFake 測試集時，無論是哪種範式皆出現了顯著的防禦力衰退，顯示未知環境噪音、通道失配與強烈有損壓縮對深偽特徵的掩蔽效應十分嚴重。未來的系統應整合專門針對語音還原與雜訊抑制的前處理網路，或在訓練階段引入對抗性生成訓練，動態生成足以混淆模型的邊界樣本，以提升真實場景下的通用防禦能力。

最後，正如 MFA-Conformer 規模擴容實驗所揭示，單純增加單一預訓練前端的模型規模無法有效突破跨域效能飽和點。參考近期在未知領域防偽任務上的前沿發展趨勢，未來預訓練—微調範式的研究重心應從「單一前端的規模堆疊」轉移至「多重自監督學習（Multi-SSL）前端的融合」。由於不同的 SSL 模型（如 WavLM、HuBERT、XLSR 等）對聲學特徵的偏好與注意力機制各異，設計有效的聚合網路將多種前端特徵截長補短地結合，將有望更全面地覆蓋並捕捉各式未知深偽生成痕跡，從而突破當前單一架構的效能天花板。
